\documentclass[10pt,letterpaper,sans]{moderncv}
\moderncvstyle{banking}
\moderncvcolor{blue}

\usepackage[scale=0.84]{geometry}
\nopagenumbers{}
\setlength{\hintscolumnwidth}{2.35cm}

\name{Krishnaa}{Vadivel}
\title{Computational Physics | Distributed Computing | Machine Learning}
\email{srikrishnaa.careers@gmail.com}
\phone[mobile]{516-853-5663}
\social[linkedin]{sri-krishnaa-ece-phd}
\extrainfo{\href{https://krishnaa423.github.io/personal-website/}{Website} \textbar{} \href{https://github.com/krishnaa423}{GitHub: krishnaa423} \textbar{} \href{https://leetcode.com/u/krishnaa42342}{LeetCode: krishnaa42342} \textbar{} \href{https://hub.docker.com/u/krishnaa42342}{Docker Hub: krishnaa42342}}

\begin{document}
\makecvtitle

\section{Profile}
\cvitem{}{Computational physics PhD researcher with training across condensed matter, numerical methods, distributed scientific software, and machine-learning-enabled materials modeling. Current work centers on first-principles calculations of self-trapped excitons and excitonic polarons, many-body theory for electron-phonon coupled systems, and scalable implementations for national-lab HPC platforms.}

\section{Research Experience}
\cventry{2021--present}{Research Assistant / PhD Researcher}{Yale University, Electrical Engineering}{}{}{\begin{itemize}%
\item Developed first-principles research workflows for self-trapped excitons, excitonic polarons, and related excited-state phenomena in materials, linking DFT, DFPT, GW/BSE-style methods, and effective Hamiltonian viewpoints.
\item Studied many-body and field-theoretic formulations of coupled electron, exciton, and phonon systems, including coherent-state and path-integral approaches for time-dependent and self-localized states.
\item Built distributed scientific computing pipelines in Python and C++ with MPI, OpenMP, PETSc, and SLEPc for large linear-algebra and simulation workloads.
\item Ported and optimized computational workloads for accelerator-enabled HPC environments using CUDA and ROCm, with experience running on Frontier, Perlmutter, and other leadership-class clusters.
\item Used PyTorch, PyTorch Geometric, and equivariant graph neural networks to accelerate self-trapped-exciton calculations and related scientific-computing workflows.
\item Packaged reproducible software in custom Docker containers and debugged distributed Linux applications with Bash and command-line tooling across HPC systems.
\item Co-authored a 2025 \textit{Journal of the American Chemical Society} paper and presented APS talks in 2024, 2025, and 2026 on self-trapped exciton formation and related excited-state materials physics.
\end{itemize}}

\section{Professional Experience}
\cventry{2019--2021}{Software and Embedded Systems Engineer}{Probe Technologies Inc. / Weatherford Wireline Services}{}{}{\begin{itemize}%
\item Built CUDA-based data-processing pipelines and custom GPU-enabled engineering tooling for large acoustic waveform datasets generated by downhole logging tools.
\item Developed GUI applications with C\#, ASP.NET server-side services, and Microsoft SQL Server workflows to automate remote calibration of temperature and pressure sensors in real time, improving production throughput.
\item Built low-level debugging tools and engineering software for deployed instrumentation, embedded-adjacent workflows, and server-side programs.
\item Applied FEniCS-based finite-element analysis for frequency analysis of steel pipes in well environments to support in-house physics and engineering studies.
\end{itemize}}
\cventry{2018--2019}{Photonics Technical Writer}{FindLight}{}{}{\begin{itemize}%
\item Produced technical articles and product-facing photonics content covering lasers, optics, and optoelectronic components for a specialized engineering readership.
\end{itemize}}

\section{Selected Projects}
\cventry{}{Distributed First-Principles Exciton-Phonon Simulation}{}{}{}{\begin{itemize}%
\item Developed computational workflows for first-principles excited-state and exciton-phonon calculations, with an emphasis on self-trapped excitons, scalable many-body theory formulations, and reusable HPC pipelines.
\end{itemize}}
\cventry{}{Equivariant Graph Neural Networks for Molecular Systems}{}{}{}{\begin{itemize}%
\item Built graph-based and symmetry-aware ML models with PyTorch and PyTorch Geometric for molecular-property prediction and explored their use as accelerators for scientific workflows.
\end{itemize}}
\cventry{}{Agentic Scientific Input Generation}{}{}{}{\begin{itemize}%
\item Built retrieval-driven tooling that reads software documentation and research papers to generate structured computational-science inputs and starter workflows for physics and chemistry codes.
\end{itemize}}
\cventry{}{Integrated Photonics and FDTD Modeling}{}{}{}{\begin{itemize}%
\item Designed photonic structures including nanophotonic ring-resonator band-pass filters in telecom bands with Meep-based FDTD workflows.
\end{itemize}}
\cventry{}{FEniCS-Based Finite-Element Analysis}{}{}{}{\begin{itemize}%
\item Built FEniCS-based finite-element workflows for frequency analysis of steel pipes in well environments to support physics-driven engineering studies.
\end{itemize}}

\section{Education}
\cventry{2021--present}{PhD, Electrical Engineering}{Yale University}{}{}{Ongoing doctoral research in computational materials, many-body theory, and scientific computing.}
\cventry{2023}{MPhil, Electrical Engineering}{Yale University}{}{}{}
\cventry{2023}{MS, Electrical Engineering}{Yale University}{}{}{}
\cventry{2017--2018}{MEng, Electrical and Computer Engineering}{Cornell University}{}{}{Included work in stochastic computing, device-oriented EE topics, and advanced design coursework.}
\cventry{2013--2017}{BS, Electrical and Computer Engineering}{Cornell University}{}{}{Included embedded systems, robotics, and core mathematics, physics, and computing foundations.}

\section{Selected Publications}
\cvitem{2025}{de Paula, A. M., Li, S., Hou, B., Vadivel, S., Teles-Ferreira, D. C., Iudica, A., Kabacinski, P., Hosseini, H., McArthur, J., Cerullo, G., et al. (2025). Time-domain observation of ultrafast self-trapped exciton formation in lead-free double halide perovskites. \textit{Journal of the American Chemical Society, 147}(32), 28923--28931.}

\section{Selected Conference Presentations}
\cvitem{2026}{Vadivel, S., Grande, R. D., Strubbe, D. A., \& Qiu, D. Y. (2026). First-principles study of exciton-polarons and self-trapped excitons. \textit{APS Global Physics Summit 2026}.}
\cvitem{2025}{Vadivel, S., Qiu, D. Y., Grande, R. D., \& Strubbe, D. A. (2025). First-principles study of self-trapped exciton formation in double perovskites. \textit{APS Global Physics Summit 2025}.}
\cvitem{2024}{Vadivel, S., \& Qiu, D. (2024). First-principles study of self-trapped exciton formation in double perovskites using excited state forces. \textit{APS March Meeting Abstracts 2024}, F01--007.}

\section{Selected Coursework}
\cvitem{Mathematics}{Differential Geometry; Abstract Algebra; Honors Mathematical Analysis II}
\cvitem{Computer Science}{Theoretical Computer Science; Probabilistic Machine Learning}
\cvitem{Physics}{Graduate Quantum Mechanics I--II; Solid State Physics I--II; Many-Body Quantum Physics}
\cvitem{Electrical Engineering}{Semiconductor Physics; Thin Film Science; Lasers and Optoelectronics; Integrated Photonics; VLSI Design; Analog VLSI Design; Mathematics of Signals and Systems; Advanced Embedded Systems}

\section{Technical Skills}
\cvitem{Languages}{Python, C, C++, CUDA, JavaScript, Bash, PowerShell, SQL, C\#}
\cvitem{Scientific Computing}{MPI, OpenMP, PETSc, SLEPc, NumPy, pandas, SciPy, SymPy, matplotlib, mpi4py, petsc4py, slepc4py, FEniCS}
\cvitem{Machine Learning}{PyTorch, PyTorch Geometric, scikit-learn, equivariant graph neural networks}
\cvitem{Systems}{Linux command line, macOS, Windows, CMake, Docker, Docker Compose, Kubernetes, gcloud, ASP.NET, Microsoft SQL Server}
\cvitem{Debugging / Platforms}{Distributed-application debugging, instrumentation-aware debugging tools, Frontier, Frontera, Perlmutter, and related national-lab or university HPC environments}

\end{document}
